\chapter{2003年辽宁卷}
\section{单选题}
\begin{problem}
与曲线 \( y = \frac{1}{x - 1} \) 关于原点对称的曲线为（\quad）

\choices
    {\( y = \frac{1}{1 + x} \)}
    {\( y = -\frac{1}{1 + x} \)}
    {\( y = \frac{1}{1 - x} \)}
    {\( y = -\frac{1}{1 - x} \)}
\end{problem}

\begin{answer}
A
\end{answer}

\begin{solution}
原曲线上点 \((x,y)\) 满足 \(y=\frac{1}{x-1}\). 关于原点对称后，点 \((x,y)\) 对应原曲线上的点 \((-x,-y)\)，故有
\[
-y=\frac{1}{-x-1}=-\frac{1}{x+1},
\]
即对称曲线的方程为 \(y=\frac{1}{x+1}\). 故选 A.
\end{solution}

\begin{problem}
已知 \(x \in \left(-\frac{\pi}{2}, 0\right)\)，\(\cos x = \frac{4}{5}\)，则 \(\tan 2x =\)（\quad）

\choices
    {\(\frac{7}{24}\)}
    {\(-\frac{7}{24}\)}
    {\(\frac{24}{7}\)}
    {\(-\frac{24}{7}\)}
\end{problem}

\begin{answer}
D
\end{answer}

\begin{solution}
因为 \(x\in\left(-\frac{\pi}{2},0\right)\)，所以 \(\sin x<0\). 由 \(\cos x=\frac45\) 得
\[
\sin x=-\sqrt{1-\cos^2x}=-\frac35.
\]
由此 \(\tan x=\frac{\sin x}{\cos x}=-\frac34\).
因此
\[
\tan 2x=\frac{2\tan x}{1-\tan^2x}
=\frac{-\frac32}{1-\frac{9}{16}}
=-\frac{24}{7}.
\]
故选 D.
\end{solution}

\begin{problem}
\(\frac{1 - \sqrt{3}\i}{(\sqrt{3} + \i)^2} =\)（\quad）

\choices
    {\(\frac{1}{4} + \frac{\sqrt{3}}{4}\i\)}
    {\(-\frac{1}{4} - \frac{\sqrt{3}}{4}\i\)}
    {\(\frac{1}{2} + \frac{\sqrt{3}}{2}\i\)}
    {\(-\frac{1}{2} - \frac{\sqrt{3}}{2}\i\)}
\end{problem}

\begin{answer}
B
\end{answer}

\begin{solution}
先计算分母
\[
(\sqrt{3}+\i)^2=2+2\sqrt{3}\i=2\left(1+\sqrt{3}\i\right).
\]
于是
\[
\frac{1-\sqrt{3}\i}{(\sqrt{3}+\i)^2}
=\frac{(1-\sqrt{3}\i)^2}{2\left(1+\sqrt{3}\i\right)\left(1-\sqrt{3}\i\right)}
=\frac{-2-2\sqrt{3}\i}{8}
=-\frac14-\frac{\sqrt{3}}4\i.
\]
故选 B.
\end{solution}

\begin{problem}
已知四边形 \(ABCD\) 是菱形，点 \(P\) 在对角线 \(AC\) 上（不包括端点 \(A\)，\(C\)），则 \(\overrightarrow{AP} =\)（\quad）

\choices
    {\(\lambda (\overrightarrow{AB} + \overrightarrow{AD})\)，\(\lambda \in (0,1)\)}
    {\(\lambda (\overrightarrow{AB} + \overrightarrow{BC})\)，\(\lambda \in \left(0,\frac{\sqrt{2}}{2}\right)\)}
    {\(\lambda (\overrightarrow{AB} - \overrightarrow{AD})\)，\(\lambda \in (0,1)\)}
    {\(\lambda (\overrightarrow{AB} - \overrightarrow{BC})\)，\(\lambda \in \left(0,\frac{\sqrt{2}}{2}\right)\)}
\end{problem}

\begin{answer}
A
\end{answer}

\begin{solution}
因为 \(ABCD\) 是菱形，所以也是平行四边形，从而
\[
\overrightarrow{AC}=\overrightarrow{AB}+\overrightarrow{AD}.
\]
点 \(P\) 在对角线 \(AC\) 的内部，故存在 \(\lambda\in(0,1)\)，使得
\[
\overrightarrow{AP}=\lambda\overrightarrow{AC}
=\lambda\left(\overrightarrow{AB}+\overrightarrow{AD}\right).
\]
故选 A.
\end{solution}

\begin{problem}
设函数 \( f(x) = \begin{cases}
2^{-x} - 1, & x \leqslant 0,\\
x^{\frac{1}{2}}, & x > 0,
\end{cases} \) 若 \( f(x_0) > 1 \)，则 \( x_0 \) 的取值范围是（\quad）

\choices
    {\((-1, 1)\)}
    {\((-1, +\infty)\)}
    {\((-\infty, -2) \cup (0, +\infty)\)}
    {\((-\infty, -1) \cup (1, +\infty)\)}
\end{problem}

\begin{answer}
D
\end{answer}

\begin{solution}
当 \(x\leqslant0\) 时，
\[
f(x)>1\Longleftrightarrow 2^{-x}-1>1
\Longleftrightarrow 2^{-x}>2
\Longleftrightarrow x<-1.
\]
当 \(x>0\) 时，
\[
f(x)>1\Longleftrightarrow \sqrt{x}>1
\Longleftrightarrow x>1.
\]
所以 \(x_0\in(-\infty,-1)\cup(1,+\infty)\). 故选 D.
\end{solution}

\begin{problem}
等差数列 \(\{a_n\}\) 中，已知 \(a_1 = \frac{1}{3}\)，\(a_2 + a_5 = 4\)，\(a_n = 33\)，则 \(n\) 为（\quad）

\choices
    {48}
    {49}
    {50}
    {51}
\end{problem}

\begin{answer}
C
\end{answer}

\begin{solution}
设等差数列的公差为 \(d\)，则 \(a_2=\frac13+d\)，\(a_5=\frac13+4d\).
由 \(a_2+a_5=4\) 得
\[
\frac23+5d=4.
\]
所以 \(d=\frac23\).
因此
\[
a_n=a_1+(n-1)d=\frac13+\frac23(n-1)=\frac{2n-1}{3}.
\]
由 \(a_n=33\) 得 \(2n-1=99\)，所以 \(n=50\). 故选 C.
\end{solution}

\begin{problem}
函数 \(y = \ln \frac{x + 1}{x - 1}\)，\(x \in (1, +\infty)\) 的反函数为（\quad）

\choices
    {\( y = \frac{\e^{x} - 1}{\e^{x} + 1}\)，\(x \in (0, +\infty)\)}
    {\(y = \frac{\e^{x} + 1}{\e^{x} - 1}\)，\(x \in (0, +\infty)\)}
    {\(y = \frac{\e^{x} - 1}{\e^{x} + 1}\)，\(x \in (-\infty, 0)\)}
    {\(y = \frac{\e^{x} + 1}{\e^{x} - 1}\)，\(x \in (-\infty, 0)\)}
\end{problem}

\begin{answer}
B
\end{answer}

\begin{solution}
设原函数的函数值为 \(y\)，则 \(y=\ln\frac{x+1}{x-1}\)，且 \(\e^y=\frac{x+1}{x-1}\).
因为原函数的定义域为 \((1,+\infty)\)，其值域为 \((0,+\infty)\). 解关于 \(x\) 的方程得
\[
\e^y(x-1)=x+1.
\]
解得 \(x=\frac{\e^y+1}{\e^y-1}\).
交换自变量和因变量，反函数为
\[
y=\frac{\e^x+1}{\e^x-1}
\]
其定义域为 \(x\in(0,+\infty)\).
故选 B.
\end{solution}

\begin{problem}
棱长为 \(a\) 的正方体中，连结相邻面的中心，以这些线段为棱的八面体的体积为（\quad）

\choices
    {\(\frac{a^3}{3}\)}
    {\(\frac{a^3}{4}\)}
    {\(\frac{a^3}{6}\)}
    {\(\frac{a^3}{12}\)}
\end{problem}

\begin{answer}
C
\end{answer}

\begin{solution}
连接六个面的中心后，所得八面体可分为两个全等的四棱锥. 每个四棱锥的底面是一个菱形，其两条对角线都为 \(a\)，所以底面积为 \(\frac{a^2}{2}\)，高为 \(\frac a2\). 因而八面体体积为
\[
2\times\frac13\times\frac{a^2}{2}\times\frac a2=\frac{a^3}{6}.
\]
故选 C.
\end{solution}

\begin{problem}
设 \(a > 0\)，\(f(x) = ax^2 + bx + c\). 曲线 \(y = f(x)\) 在点 \(P(x_0, f(x_0))\) 处切线的倾斜角的取值范围为 \(\left[0, \frac{\pi}{4}\right]\)，则 \(P\) 到曲线 \(y = f(x)\) 对称轴距离的取值范围为（\quad）

\choices
    {\(\left[0, \frac{1}{a}\right]\)}
    {\(\left[0, \frac{1}{2a}\right]\)}
    {\(\left[0, \left|\frac{b}{2a}\right|\right]\)}
    {\(\left[0, \left|\frac{b - 1}{2a}\right|\right]\)}
\end{problem}

\begin{answer}
B
\end{answer}

\begin{solution}
抛物线在点 \(P\) 处切线的斜率为 \(f'(x_0)=2ax_0+b\). 因为切线倾斜角属于 \(\left[0,\frac\pi4\right]\)，所以
\[
0\leqslant2ax_0+b=\tan\theta\leqslant1.
\]
抛物线的对称轴为 \(x=-\frac{b}{2a}\)，故点 \(P\) 到对称轴的距离为
\[
\left|x_0+\frac{b}{2a}\right|=\frac{|2ax_0+b|}{2a}=\frac{\tan\theta}{2a}.
\]
因此该距离的取值范围为 \(\left[0,\frac{1}{2a}\right]\). 故选 B.
\end{solution}

\begin{problem}
已知双曲线中心在原点且一个焦点为 \(F(\sqrt{7}, 0)\)，直线 \(y = x - 1\) 与其相交于 \(M\)，\(N\) 两点，\(MN\) 中点的横坐标为 \(-\frac{2}{3}\)，则此双曲线的方程是（\quad）

\choices
    {\(\frac{x^2}{3} - \frac{y^2}{4} = 1\)}
    {\(\frac{x^2}{4} - \frac{y^2}{3} = 1\)}
    {\(\frac{x^2}{5} - \frac{y^2}{2} = 1\)}
    {\(\frac{x^2}{2} - \frac{y^2}{5} = 1\)}
\end{problem}

\begin{answer}
D
\end{answer}

\begin{solution}
设双曲线方程为
\[
\frac{x^2}{p}-\frac{y^2}{q}=1.
\]
其中 \(p>0\)，\(q>0\).
由焦点为 \(F(\sqrt7,0)\) 得
\[
p+q=7.
\]
将直线 \(y=x-1\) 代入双曲线方程，得交点横坐标满足
\[
(q-p)x^2+2px-p-pq=0.
\]
两交点横坐标的平均值就是弦 \(MN\) 中点的横坐标，因此
由中点横坐标条件得 \(-\frac{p}{q-p}=-\frac23\)，即 \(3p=2(q-p)\).
联立 \(p+q=7\)，解得 \(p=2\)，\(q=5\). 所以双曲线方程为
\[
\frac{x^2}{2}-\frac{y^2}{5}=1.
\]
故选 D.
\end{solution}

\begin{problem}
已知长方形的四个顶点 \(A(0,0)\)，\(B(2,0)\)，\(C(2,1)\) 和 \(D(0,1)\)，一质点从 \(AB\) 的中点 \(P_0\) 沿与 \(AB\) 的夹角 \(\theta\) 的方向射到 \(BC\) 上的点 \(P_1\) 后，依次反射到 \(CD\)，\(DA\) 和 \(AB\) 上的点 \(P_2\)，\(P_3\) 和 \(P_4\)（入射角等于反射角），设 \(P_4\) 的坐标为 \((x_4,0)\)，若 \(1 < x_4 < 2\)，则 \(\tan \theta\) 的取值范围是（\quad）

\choices
    {\(\left(\frac{1}{3},1\right)\)}
    {\(\left(\frac{1}{3},\frac{2}{3}\right)\)}
    {\(\left(\frac{2}{5},\frac{1}{2}\right)\)}
    {\(\left(\frac{2}{5},\frac{2}{3}\right)\)}
\end{problem}

\begin{answer}
C
\end{answer}

\begin{solution}
设 \(t=\tan\theta\). 质点从 \(P_0(1,0)\) 射向右侧，先到达 \(BC\) 上的
\[
P_1=(2,t).
\]
其中 \(0<t<1\).
反射到上边 \(CD\) 时，水平位移与竖直位移之比为 \(1:t\)，所以
\[
P_2=\left(2-\frac{1-t}{t},1\right)=\left(3-\frac1t,1\right).
\]
再反射到左边 \(DA\)，得到
\[
P_3=\left(0,1-t\left(3-\frac1t\right)\right)=(0,2-3t).
\]
最后反射到 \(AB\)，故
\[
x_4=\frac{2-3t}{t}=\frac2t-3.
\]
由各次反射点均落在线段内部，至少有 \(\frac13<t<\frac23\). 题设又给出
\[
1<\frac2t-3<2.
\]
因为 \(t>0\)，解得 \(\frac25<t<\frac12\). 故选 C.
\end{solution}

\begin{problem}
一个四面体的所有棱长都为 \(\sqrt{2}\)，四个顶点在同一球面上，则此球的表面积为（\quad）

\choices
    {\(3\pi\)}
    {\(4\pi\)}
    {\(3\sqrt{3}\pi\)}
    {\(6\pi\)}
\end{problem}

\begin{answer}
A
\end{answer}

\begin{solution}
可将四面体的四个顶点取为
\[
\left(\frac12,\frac12,\frac12\right),
\left(\frac12,-\frac12,-\frac12\right),
\left(-\frac12,\frac12,-\frac12\right),
\left(-\frac12,-\frac12,\frac12\right).
\]
任意两点恰有两个坐标相差 \(1\)，所以各棱长均为 \(\sqrt2\)，且四点到原点的距离均为
\[
R=\sqrt{\frac14+\frac14+\frac14}=\frac{\sqrt3}{2}.
\]
故外接球表面积为
\[
4\pi R^2=4\pi\times\frac34=3\pi.
\]
故选 A.
\end{solution}

\section{填空题}
\begin{problem}
\(\left(x^{2} - \frac{1}{2x}\right)^{9}\) 的展开式中 \(x^{9}\) 系数是\fillinblank{}.
\end{problem}

\begin{answer}
\(-\frac{21}{2}\)
\end{answer}

\begin{solution}
展开式通项中，若从第二项 \(-\frac{1}{2x}\) 取 \(k\) 次，则该项为
\[
\mathrm C_9^k\left(x^2\right)^{9-k}\left(-\frac{1}{2x}\right)^k,
\]
其中 \(x\) 的指数为 \(2(9-k)-k=18-3k\). 要得到 \(x^9\)，必须有
\[
18-3k=9.
\]
所以 \(k=3\).
所以所求系数为
\[
\mathrm C_9^3\left(-\frac12\right)^3=84\times\left(-\frac18\right)=-\frac{21}{2}.
\]
\end{solution}

\begin{problem}
某公司生产三种型号的轿车，产量分别为 1200 辆，6000 辆和 2000 辆. 为检验该公司的产品质量，现用分层抽样的方法抽取 46 辆进行检验，这三种型号的轿车依次应抽取\fillinblank{}，\fillinblank{}，\fillinblank{}辆.
\end{problem}

\begin{answer}
\(6,30,10\)
\end{answer}

\begin{solution}
三种型号轿车的总产量为 \(1200+6000+2000=9200\) 辆. 分层抽样中各层的抽样比例相同，故抽样比例为
\[
\frac{46}{9200}=\frac1{200}.
\]
因此三种型号依次抽取 \(1200\times\frac1{200}=6\)，\(6000\times\frac1{200}=30\)，\(2000\times\frac1{200}=10\) 辆.
\end{solution}

\begin{problem}
某城市在中心广场建造一个花圃，花圃分为 6 个部分（如图）. 现要栽种 4 种不同颜色的花，每部分栽种一种且相邻部分不能栽种同样颜色的花，不同的栽种方法有\fillinblank{}种. （以数字作答）

\begin{center}
\bitmapfigure[max width=6.0cm]{img/2003/liaoning/q15_fig1.png}
\end{center}
\end{problem}

\begin{answer}
120
\end{answer}

\begin{solution}
由图可知，中心部分 1 与外部五个部分均相邻，外部五个部分依次相邻并构成一个五边形. 先给中心部分 1 选颜色，有 \(4\) 种选法；外部各部分不能使用中心部分的颜色，因此外部五边形相当于用其余 \(3\) 种颜色进行相邻异色染色.

固定部分 2 的颜色，有 \(3\) 种选法. 从部分 2 到部分 6 依次给部分 3，4，5，6 染色，部分 6 还必须与部分 2 异色. 部分 3，4，5 共有 \(2^3=8\) 种染法，其中部分 5 与部分 2 同色的有 \(2\) 种，此时部分 6 有 \(2\) 种选法；部分 5 与部分 2 异色的有 \(6\) 种，此时部分 6 只有 \(1\) 种选法. 因而外部五边形有
\[
3\left(2\times2+6\times1\right)=30
\]
种染法. 总方法数为
\[
4\times30=120.
\]
\end{solution}

\begin{problem}
对于四面体 \(ABCD\)，给出下列四个命题：

\begin{circlelist}
\item 若 \(AB = AC\)，\(BD = CD\)，则 \(BC \perp AD\)；
\item 若 \(AB = CD\)，\(AC = BD\)，则 \(BC \perp AD\)；
\item 若 \(AB \perp AC\)，\(BD \perp CD\)，则 \(BC \perp AD\)；
\item 若 \(AB \perp CD\)，\(AC \perp BD\)，则 \(BC \perp AD\).
\end{circlelist}

其中真命题的序号是\fillinblank{}. （写出所有真命题的序号）
\end{problem}

\begin{answer}
\(\circled{1}\circled{4}\)
\end{answer}

\begin{solution}
令 \(\bs{u}=\overrightarrow{AB}\)，\(\bs{v}=\overrightarrow{AC}\)，\(\bs{w}=\overrightarrow{AD}\).
对于命题 \(\circled{1}\)，由 \(AB=AC\) 得
\[
(\bs{u}-\bs{v})\cdot(\bs{u}+\bs{v})=0.
\]
由 \(BD=CD\) 得
\[
|\bs{u}-\bs{w}|^2=|\bs{v}-\bs{w}|^2.
\]
展开得 \((\bs{u}-\bs{v})\cdot(\bs{u}+\bs{v}-2\bs{w})=0\).
两式相减得 \((\bs{u}-\bs{v})\cdot\bs{w}=0\). 由于 \(\overrightarrow{BC}=\bs{v}-\bs{u}\)，所以 \(BC\perp AD\)，命题 \(\circled{1}\) 为真.

命题 \(\circled{2}\) 不一定成立. 取
\(A=(0,0,0)\)，\(B=(1,0,0)\)，\(C=(0,1,1)\)，\(D=(0,0,1)\).
此时 \(AB=CD=1\)，\(AC=BD=\sqrt2\)，但
\[
\overrightarrow{BC}\cdot\overrightarrow{AD}=(-1,1,1)\cdot(0,0,1)=1\neq0,
\]
所以命题 \(\circled{2}\) 为假.

命题 \(\circled{3}\) 也不一定成立. 取
\(A=(0,0,0)\)，\(B=(1,0,0)\)，\(C=(0,1,0)\)，\(D=\left(1,\frac12,\frac12\right)\).
有 \(\overrightarrow{AB}\cdot\overrightarrow{AC}=0\)，且
\[
\overrightarrow{BD}\cdot\overrightarrow{CD}
=\left(0,\frac12,\frac12\right)\cdot\left(1,-\frac12,\frac12\right)=0,
\]
但
\[
\overrightarrow{BC}\cdot\overrightarrow{AD}=(-1,1,0)\cdot\left(1,\frac12,\frac12\right)=-\frac12\neq0,
\]
所以命题 \(\circled{3}\) 为假.

对于命题 \(\circled{4}\)，\(AB\perp CD\) 和 \(AC\perp BD\) 分别给出
\[
\bs{u}\cdot(\bs{v}-\bs{w})=0.
\]
并且 \(\bs{v}\cdot(\bs{u}-\bs{w})=0\).
因此 \(\bs{u}\cdot\bs{w}=\bs{u}\cdot\bs{v}=\bs{v}\cdot\bs{w}\)，从而
\[
\overrightarrow{BC}\cdot\overrightarrow{AD}=(\bs{v}-\bs{u})\cdot\bs{w}=0.
\]
所以命题 \(\circled{4}\) 为真，真命题的序号为 \(\circled{1}\circled{4}\).
\end{solution}

\section{解答题}
\begin{problem}
已知正四棱柱 \(ABCD - A_{1}B_{1}C_{1}D_{1}\)，\(AB = 1\)，\(AA_{1} = 2\)，\(E\) 为 \(CC_{1}\) 中点，\(F\) 为 \(BD_{1}\) 中点.

\begin{enumerate}
\item 证明： \(EF\) 为 \(BD_{1}\) 与 \(CC_{1}\) 的公垂线；
\item 求点 \(D_{1}\) 到面 \(BDE\) 的距离.
\end{enumerate}

\bitmapfigure[max width=4.8cm]{img/2003/liaoning/q17_fig1.png}
\end{problem}

\begin{solution}
取 \(A\) 为原点，令 \(AB\)，\(AD\)，\(AA_{1}\) 的方向分别为 \(x\)，\(y\)，\(z\) 轴，则坐标为 \(A(0,0,0)\)，\(B(1,0,0)\)，\(D(0,1,0)\)，\(C(1,1,0)\)，\(D_{1}(0,1,2)\)，\(C_{1}(1,1,2)\)，\(E(1,1,1)\)，\(F\left(\frac{1}{2},\frac{1}{2},1\right)\).

\begin{enumerate}
\item 由坐标可得 \(\overrightarrow{EF}=\left(-\frac{1}{2},-\frac{1}{2},0\right)\)，\(\overrightarrow{BD_{1}}=(-1,1,2)\)，\(\overrightarrow{CC_{1}}=(0,0,2)\). 因此 \(\overrightarrow{EF}\cdot\overrightarrow{BD_{1}}=0\)，且 \(\overrightarrow{EF}\cdot\overrightarrow{CC_{1}}=0\).
又 \(F\in BD_{1}\)，\(E\in CC_{1}\)，所以 \(EF\) 同时垂直于 \(BD_{1}\) 和 \(CC_{1}\)，故 \(EF\) 是 \(BD_{1}\) 与 \(CC_{1}\) 的公垂线.

\item 由 \(\overrightarrow{BD}=(-1,1,0)\)，\(\overrightarrow{BE}=(0,1,1)\)，
可取平面 \(BDE\) 的一个法向量为
\[
\bs{n}=\overrightarrow{BD}\times\overrightarrow{BE}=(1,1,-1).
\]
所以平面 \(BDE\) 的方程为 \(x+y-z-1=0\).
点 \(D_{1}(0,1,2)\) 到平面 \(BDE\) 的距离为
\[
\frac{|0+1-2-1|}{\sqrt{1^{2}+1^{2}+(-1)^{2}}}
=\frac{2}{\sqrt{3}}=\frac{2\sqrt{3}}{3}.
\]
\end{enumerate}
\end{solution}

\begin{problem}
已知函数 \( f(x) = \sin (\omega x + \varphi) \)（\(\omega > 0\)，\(0 \leqslant \varphi \leqslant \pi\)）是 \(\R\) 上的偶函数，其图象关于点 \( M\left(\frac{3\pi}{4}, 0\right) \) 对称，且在区间 \(\left[0, \frac{\pi}{2}\right]\) 上是单调函数. 求 \(\omega\) 和 \(\varphi\) 的值.
\end{problem}

\begin{solution}
由 \(f(x)\) 是偶函数，对任意 \(x\in\R\) 有
\[
0=f(x)-f(-x)=\sin(\varphi+\omega x)-\sin(\varphi-\omega x)
=2\cos\varphi\sin(\omega x).
\]
由于 \(\omega>0\)，上式对任意 \(x\) 成立只能说明 \(\cos\varphi=0\). 结合 \(0\leqslant\varphi\leqslant\pi\)，得 \(\varphi=\frac{\pi}{2}\)，且 \(f(x)=\cos(\omega x)\).

图象关于 \(M\left(\frac{3\pi}{4},0\right)\) 对称，等价于对任意 \(t\in\R\) 有
\[
f\left(\frac{3\pi}{4}+t\right)+f\left(\frac{3\pi}{4}-t\right)=0.
\]
代入 \(f(x)=\cos(\omega x)\)，得
\[
2\cos\frac{3\pi\omega}{4}\cos(\omega t)=0.
\]
因此 \(\cos\frac{3\pi\omega}{4}=0\)，所以存在 \(k\in\Z\)，使得 \(\frac{3\pi\omega}{4}=\frac{(2k+1)\pi}{2}\).
即
\[
\omega=\frac{4k+2}{3}.
\]

在区间 \(\left[0,\frac{\pi}{2}\right]\) 上，函数 \(\cos(\omega x)\) 的自变量区间为 \(\left[0,\frac{\omega\pi}{2}\right]\). 要使其在该区间上单调，必须有 \(0<\frac{\omega\pi}{2}\leqslant\pi\)，即 \(0<\omega\leqslant2\). 否则自变量区间会越过 \(\pi\)，余弦函数先减后增而不单调. 由 \(\omega=\frac{4k+2}{3}\) 和 \(0<\omega\leqslant2\)，得 \(\omega=\frac{2}{3}\) 或 \(\omega=2\). 故 \(\varphi=\frac{\pi}{2}\)，且 \(\omega=\frac{2}{3}\) 或 \(\omega=2\).
\end{solution}

\begin{problem}
设 \(a > 0\)，求函数 \(f(x) = \sqrt{x} - \ln (x + a)\)，\(x \in (0, +\infty)\) 的单调区间.
\end{problem}

\begin{solution}
当 \(x>0\) 时，
\[
f'(x)=\frac{1}{2\sqrt{x}}-\frac{1}{x+a}
=\frac{x+a-2\sqrt{x}}{2\sqrt{x}(x+a)}
=\frac{(\sqrt{x}-1)^{2}+a-1}{2\sqrt{x}(x+a)}.
\]
分母始终为正，因此只需讨论分子符号.

当 \(a\geqslant1\) 时，
\[
(\sqrt{x}-1)^{2}+a-1\geqslant0.
\]
当 \(a=1\) 时，分子只在 \(x=1\) 处为零；当 \(a>1\) 时，分子始终为正. 因此函数在 \((0,+\infty)\) 上单调递增.

当 \(0<a<1\) 时，令 \(r=\sqrt{1-a}\)，则 \(0<r<1\)，且
\[
(\sqrt{x}-1)^{2}+a-1=(\sqrt{x}-1)^{2}-r^{2}
=\bigl(\sqrt{x}-1-r\bigr)\bigl(\sqrt{x}-1+r\bigr).
\]
两个零点分别为 \(x_{1}=(1-r)^{2}=\bigl(1-\sqrt{1-a}\bigr)^{2}\)，\(x_{2}=(1+r)^{2}=\bigl(1+\sqrt{1-a}\bigr)^{2}\).
于是
当 \(x\in(0,x_{1})\cup(x_{2},+\infty)\) 时，\(f'(x)>0\)；当 \(x\in(x_{1},x_{2})\) 时，\(f'(x)<0\).
所以当 \(0<a<1\) 时，函数在 \((0,x_{1})\) 和 \((x_{2},+\infty)\) 上单调递增，在 \((x_{1},x_{2})\) 上单调递减.
\end{solution}

\begin{problem}
\(A\)，\(B\) 两个代表队进行乒乓球对抗赛，每队三名队员，\(A\) 队队员是 \(A_1\)，\(A_2\)，\(A_3\)，\(B\) 队队员是 \(B_1\)，\(B_2\)，\(B_3\)，按以往多次比赛的统计，对阵队员之间胜负概率如下：

\begin{center}
\begin{tabular}{|c|c|c|}
\hline
对阵队员 & \(A\) 队队员胜的概率 & \(A\) 队队员负的概率 \\
\hline
\(A_1\) 对 \(B_1\) & \(\frac{2}{3}\) & \(\frac{1}{3}\) \\
\hline
\(A_2\) 对 \(B_2\) & \(\frac{2}{5}\) & \(\frac{3}{5}\) \\
\hline
\(A_3\) 对 \(B_3\) & \(\frac{2}{5}\) & \(\frac{3}{5}\) \\
\hline
\end{tabular}
\end{center}

现按表中对阵方式出场，每场胜队得 1 分，负队得 0 分，设 \(A\) 队，\(B\) 队最后所得总分分别为 \(\xi\)，\(\eta\).

\begin{enumerate}
\item 求 \(\xi\)，\(\eta\) 的概率分布；
\item 求 \(E(\xi)\)，\(E(\eta)\).
\end{enumerate}
\end{problem}

\begin{solution}
设三场比赛中 \(A\) 队分别获胜的事件相互独立. 由每场比赛不是 \(A\) 队胜就是 \(B\) 队胜，有
\[
\xi+\eta=3.
\]
逐项计算 \(A\) 队总得分的概率：
\[
P(\xi=3)=\frac{2}{3}\cdot\frac{2}{5}\cdot\frac{2}{5}=\frac{8}{75},
\]
\[
P(\xi=2)=\frac{2}{3}\cdot\frac{2}{5}\cdot\frac{3}{5}
+\frac{2}{3}\cdot\frac{3}{5}\cdot\frac{2}{5}
+\frac{1}{3}\cdot\frac{2}{5}\cdot\frac{2}{5}
=\frac{28}{75},
\]
\[
P(\xi=1)=\frac{1}{3}\cdot\frac{2}{5}\cdot\frac{3}{5}
+\frac{1}{3}\cdot\frac{3}{5}\cdot\frac{2}{5}
+\frac{2}{3}\cdot\frac{3}{5}\cdot\frac{3}{5}
=\frac{2}{5},
\]
\[
P(\xi=0)=\frac{1}{3}\cdot\frac{3}{5}\cdot\frac{3}{5}=\frac{3}{25}.
\]
因此分布列为
\begin{center}
\begin{tabular}{c|cccc}
\hline
\(\xi\) & 0 & 1 & 2 & 3 \\
\hline
\(P\) & \(\frac{3}{25}\) & \(\frac{2}{5}\) & \(\frac{28}{75}\) & \(\frac{8}{75}\) \\
\hline
\(\eta\) & 3 & 2 & 1 & 0 \\
\hline
\(P\) & \(\frac{3}{25}\) & \(\frac{2}{5}\) & \(\frac{28}{75}\) & \(\frac{8}{75}\) \\
\hline
\end{tabular}
\end{center}

数学期望为
\[
E(\xi)=0\cdot\frac{3}{25}+1\cdot\frac{2}{5}+2\cdot\frac{28}{75}+3\cdot\frac{8}{75}
=\frac{22}{15}.
\]
又因为 \(\eta=3-\xi\)，所以
\[
E(\eta)=3-E(\xi)=3-\frac{22}{15}=\frac{23}{15}.
\]
\end{solution}

\begin{problem}
设 \(a_0\) 为常数，且 \(a_n = 3^{n-1} - 2a_{n-1}\)（\(n \in \N^*\)）.

\begin{enumerate}
\item 证明对任意 \(n \geqslant 1\)，\(a_n = \frac{1}{5} \bigl[3^n + (-1)^{n-1} \cdot 2^n\bigr] + (-1)^n \cdot 2^n a_0\)；
\item 假设对任意 \(n \geqslant 1\) 有 \(a_n > a_{n-1}\)，求 \(a_0\) 的取值范围.
\end{enumerate}
\end{problem}

\begin{solution}
\begin{enumerate}
\item 当 \(n=1\) 时，由递推关系
\[
a_{1}=1-2a_{0}
=\frac{1}{5}\left[3+(-1)^{0}\cdot2\right]+(-1)^{1}\cdot2a_{0},
\]
结论成立.

假设当 \(n=k\)（\(k\geqslant1\)）时结论成立，即
\[
a_{k}=\frac{1}{5}\left[3^{k}+(-1)^{k-1}\cdot2^{k}\right]+(-1)^{k}\cdot2^{k}a_{0}.
\]
代入归纳假设，得
\[
a_{k+1}=3^{k}-2a_{k}
=3^{k}-\frac{2}{5}\left[3^{k}+(-1)^{k-1}\cdot2^{k}\right]-(-1)^{k}\cdot2^{k+1}a_{0}.
\]
化简得
\[
a_{k+1}=\frac{1}{5}\left[3^{k+1}+(-1)^{k}\cdot2^{k+1}\right]+(-1)^{k+1}\cdot2^{k+1}a_{0}.
\]
所以当 \(n=k+1\) 时结论也成立. 由数学归纳法，
\[
a_n=\frac{1}{5}\left[3^{n}+(-1)^{n-1}\cdot2^{n}\right]+(-1)^{n}\cdot2^{n}a_{0}
\]
对任意 \(n\geqslant1\) 成立.

\item 由上式可得
\[
5(a_n-a_{n-1})=2\cdot3^{n-1}+3(-1)^{n-1}\cdot2^{n-1}+15(-1)^{n}\cdot2^{n-1}a_{0}.
\]
由 \(n=1\) 得 \(a_{1}-a_{0}=1-3a_{0}>0\)，由 \(n=2\) 得 \(a_{2}-a_{1}=6a_{0}>0\).
故必须有
\[
0<a_{0}<\frac{1}{3}.
\]

反之，设 \(0<a_{0}<\frac{1}{3}\). 当 \(n\) 为奇数时，
\[
5(a_n-a_{n-1})=2\cdot3^{n-1}+3\cdot2^{n-1}(1-5a_{0})
>2\cdot3^{n-1}-2\cdot2^{n-1}\geqslant0.
\]
当 \(n\) 为偶数时，
\[
5(a_n-a_{n-1})=2\cdot3^{n-1}-3\cdot2^{n-1}(1-5a_{0})
>2\cdot3^{n-1}-3\cdot2^{n-1}\geqslant0.
\]
上述不等式中的第一步均为严格不等式，故 \(a_n-a_{n-1}>0\) 对任意 \(n\geqslant1\) 成立. 因此
\[
\boxed{0<a_{0}<\frac{1}{3}}.
\]
\end{enumerate}
\end{solution}

\begin{problem}
已知常数 \(a > 0\)，向量 \(\bs{c} = (0, a)\)，\(\bs{i} = (1, 0)\). 经过原点 \(O\) 以 \(\bs{c} + \lambda \bs{i}\) 为方向向量的直线与经过定点 \(A(0, a)\) 以 \(\bs{i} - 2\lambda \bs{c}\) 为方向向量的直线相交于 \(P\)，其中 \(\lambda \in \R\). 试问：是否存在两个定点 \(E\)，\(F\)，使得 \(|PE| + |PF|\) 为定值. 若存在，求出 \(E\)，\(F\) 的坐标；若不存在，说明理由.
\end{problem}

\begin{solution}
设交点 \(P\) 在两条直线上的参数分别为 \(u\)，\(v\). 由第一条直线可写成 \(P=u(\lambda,a)=\bigl(u\lambda,ua\bigr)\)，由第二条直线可写成 \(P=(0,a)+v(1,-2a\lambda)=\bigl(v,a-2a\lambda v\bigr)\).
由两式对应坐标相等，得 \(v=u\lambda\)，且
\[
ua=a-2a\lambda v=a-2au\lambda^{2}.
\]
因为 \(a>0\)，由此 \(u=\frac{1}{1+2\lambda^{2}}\)，且 \(P\left(\frac{\lambda}{1+2\lambda^{2}},\frac{a}{1+2\lambda^{2}}\right)\).
记 \(P(x,y)\)，则 \(x=\frac{\lambda}{1+2\lambda^{2}}\)，\(y=\frac{a}{1+2\lambda^{2}}\).
从而
\[
y(a-y)=\frac{a}{1+2\lambda^{2}}\cdot\frac{2a\lambda^{2}}{1+2\lambda^{2}}
=2a^{2}x^{2}.
\]
因此 \(P\) 的轨迹满足
\[
2a^{2}x^{2}+y^{2}-ay=0,
\]
即
\[
\frac{x^{2}}{1/8}+\frac{\left(y-\frac{a}{2}\right)^{2}}{a^{2}/4}=1.
\]
这是中心为 \(\left(0,\frac{a}{2}\right)\)，横半轴为 \(\alpha=\frac{1}{2\sqrt{2}}\)，纵半轴为 \(\beta=\frac{a}{2}\) 的椭圆. 由 \(y=\frac{a}{1+2\lambda^{2}}>0\) 可知原点不在实际轨迹上；反之，椭圆上任一点 \((x,y)\ne(0,0)\) 都有 \(y>0\)，取 \(\lambda=\frac{ax}{y}\)，即可由 \(x=\frac{\lambda}{1+2\lambda^{2}}\) 和 \(y=\frac{a}{1+2\lambda^{2}}\) 得到该点. 因此实际轨迹是该椭圆去掉原点后的点集，这不影响焦点距离和为定值的结论.

当 \(0<a<\frac{1}{\sqrt{2}}\) 时，\(\alpha>\beta\)，椭圆焦点在水平方向，可取
\[
E\left(\sqrt{\frac18-\frac{a^{2}}4},\frac{a}{2}\right)
\]
以及
\[
F\left(-\sqrt{\frac18-\frac{a^{2}}4},\frac{a}{2}\right)
\]
此时
\[
|PE|+|PF|=2\alpha=\frac{1}{\sqrt{2}}.
\]

当 \(a>\frac{1}{\sqrt{2}}\) 时，\(\beta>\alpha\)，椭圆焦点在竖直方向，可取
\[
E\left(0,\frac{a}{2}+\sqrt{\frac{a^{2}}4-\frac18}\right)
\]
以及
\[
F\left(0,\frac{a}{2}-\sqrt{\frac{a^{2}}4-\frac18}\right)
\]
此时
\[
|PE|+|PF|=2\beta=a.
\]

当 \(a=\frac{1}{\sqrt{2}}\) 时，轨迹为去掉原点的圆. 若存在两个不同的定点 \(E\)，\(F\)，使 \(|PE|+|PF|\) 在该轨迹上为定值，由距离函数的连续性可知，它在整个圆上也为定值. 但到两个不同定点的距离和为定值所表示的轨迹是以 \(E\)，\(F\) 为焦点的非圆椭圆，不可能是圆. 因此此时不存在符合题意的两个定点.
\end{solution}
